% ===== §3 =====
\section{The Genealogy of Optimizationism}
\label{sec:genealogy}

\noindent
The picture criticized here was not always a picture of everything, and recovering the outline of how it became one is not idle history but part of the argument. A thesis that presents itself as the timeless form of reason draws much of its authority from seeming to have no alternative and no origin, to be simply how thinking goes; and that authority is loosened, though not refuted, once the thesis is shown to have been assembled, step by step, out of choices that could have been made otherwise. The point of the genealogy is therefore not to refute optimizationism by naming its birthday, which would be the genetic fallacy, but to change its standing: to move the claim that optimization is the form of things out of the register of self-evident truth, where it defends itself by seeming to have always been there, and into the register of historical construction, where one can ask what was assumed at each step and whether the assumptions hold. Three moments mark the passage from a local operation to a picture of the world, and at each a commitment enters that the later parts will have to test.

\subsection*{From purposive action to the calculus of means}

The first moment is the isolation, within the theory of action, of a specifically instrumental rationality. Weber's typology distinguishes the rationality of ends, \emph{Wertrationalität}, oriented to a value held for its own sake, from the rationality of means, \emph{Zweckrationalität}, which takes an end as given and concerns itself with the efficient fitting of means to it \citep{weber1978economy}. The distinction is analytic, and in Weber it carries no triumph; he watched the advance of means-rationality with foreboding, as the closing of an iron cage around a disenchanted world. But the distinction did a lasting and double-edged work. By carving out a domain in which reason operates on means alone, with the end bracketed as given from outside, it made possible a rigorous science of efficient action; and it opened, at the same stroke, a fateful slippage. It is a short step, and one Weber himself did not take but the tradition after him did, from ``reason can operate on means, the end being given'' to ``reason \emph{is} the operation on means, and the end is simply given.'' Once the end is bracketed it is easily forgotten as a question, and the bracketing that was meant as a division of labor within reason hardens into a definition of reason as such. What drops out is precisely the rationality of ends, the reasoned deliberation about what is worth pursuing, which Weber had kept in view and which the instrumental turn consigns to the merely subjective, a matter of decision rather than of reason. The seed of optimizationism is here: a form of reasoning whose entire content is the maximization of the satisfaction of an end it does not itself examine.

\subsection*{The marginalist objectification of the objective}

The second moment gives the bracketed end a mathematical body and, in doing so, smuggles in the assumption on which everything later turns. In the marginalist tradition the ends of action are gathered into a single object, the utility function, and the agent is defined as its constrained maximizer; economics is recast as the science of the relation between ends and scarce means with alternative uses, and the rational agent is the one who allocates the means so as to maximize utility \citep{robbins1932nature}. Two things happen here that the criticism will need to have named. The first is a \emph{commensuration by construction}. To represent the heterogeneous ends of a life as arguments of one function returning one magnitude is already to assume that they share a scale, that gains in one can be set against losses in another and a net computed; the common scale whose existence the third impossibility will deny is not argued for in the marginalist construction but built into the form of the utility function, where it becomes invisible precisely because it is structural rather than asserted. One does not notice an assumption one has been given as the shape of the tool. The second is an \emph{objectification of the end}. What in Weber was an end held by an agent, something the agent is toward, becomes a function the agent \emph{has}, a fixed datum standing ready for the optimization to range over; and with this the endogeneity of ends, their formation and revision within the very activity that pursues them, drops silently from view. The agent of marginalism arrives at his problem with his function already written, and the writing is treated as no part of the economic act, as though preferences were furniture brought to the room rather than something that happens in the living.

\subsection*{The axiomatic universalization in decision theory}

The third moment lifts the maximizing agent clear of economics and installs him as the definition of rationality wherever there is choice at all. In the theory of expected utility, a small set of axioms on preference, of which completeness and transitivity are the load-bearing pair, is shown to entail that the agent's choices can be represented as the maximization of the expectation of a utility function; to be rational just is to choose as if maximizing such a function \citep{vonneumann1944games,savage1954foundations}. The result is a genuine achievement, elegant and, within its assumptions, unimpeachable, and the paper does not dispute the mathematics. What it marks is the quiet enormity of one axiom. Completeness requires that for any two options the agent prefer the one, prefer the other, or be indifferent, and this innocuous-seeming demand legislates that there is no pair of goods standing in any fourth relation, none between which preference is not yet settled in a way no sweetening resolves; transitivity then chains the pairwise comparisons into a single order. To accept the axioms is thus to decree, at the ground floor and before any particular choice is examined, that value is commensurable and the rational agent a maximizer on one scale, and, worse for the later argument, to make this decree the very meaning of the word ``rational.'' The consequence is that the person who declines to place two goods on a common scale, who holds that a vocation and a friendship are not rankable as better, worse, or precisely equal, is no longer voicing a substantive view about the goods that might be right; he is, by the definition now installed, irrational, his refusal a symptom rather than a position. Optimization, which entered as the efficient fitting of means to a bracketed end, has become by this third step the criterion of reason itself, binding on every chooser, in love no less than in commerce, and the completeness axiom has made the denial of commensurability unsayable within the frame, since to deny it is simply to fail the definition of the rational.

\subsection*{The reach of the genealogy, and its limit}

Set end to end, the three moments trace the elevation of a local operation into a picture of the world: the bracketing of ends that made a pure calculus of means possible, the objectification and construction-borne commensuration of ends in the utility function, and the axiomatic installation of the maximizer as the meaning of rationality as such. What the genealogy establishes is that the claim ``optimization is the form of rationality,'' and its cosmic shadow ``optimization is the form of things,'' are constructions rather than findings, and that at each step an assumption entered, commensurability above all, which could have been declined and which the paper will argue is false in the domain of intimate relation. What the genealogy does not establish, and does not pretend to, is that optimization is thereby worthless; the fallacy of dismissing a tool by its history is one the paper is careful not to commit. Optimization is a superb instrument within the domain where its four conditions hold, and the shipping fleet of the previous section is not the exception but the rule across engineering, logistics, and the allocation of genuinely fungible resources, where to refuse the maximum would be a failure of reason and not a fidelity to anything finer. The error lies in the universalization rather than in the operation, in the inference from ``this works wherever value is complete, fixed, commensurable, and owned by a stable subject'' to ``value is everywhere so, and where it seems not to be we have merely not yet done the sums.'' That inference is what the completeness axiom makes feel like a truism and what the genealogy exposes as a step that was taken and could have been refused. It is that inference, and not the mathematics of maximization, that the next part denies, condition by condition and starting from the ground.
